\documentclass[11pt]{article}

\usepackage[margin=1in]{geometry}
\usepackage{amsmath, amssymb}
\usepackage{booktabs}
\usepackage{enumitem}
\usepackage{hyperref}
\usepackage{xcolor}
\usepackage{parskip}
\usepackage{microtype}
\usepackage{tabularx}
\usepackage{array}

\hypersetup{
    colorlinks=true,
    linkcolor=black,
    urlcolor=black,
    citecolor=black
}

\title{A Beginner's Primer to the Four ML Architectures \\
       for the FXT Centrality Paper}
\author{Internal note for Mathias Labonté}
\date{\today}

\newcommand{\sqsnn}{\ensuremath{\sqrt{s_{NN}}}}
\newcommand{\pT}{\ensuremath{p_{\text{T}}}}
\newcommand{\etalab}{\ensuremath{\eta_{\text{lab}}}}

\begin{document}
\maketitle

\section*{How to read this document}

You know the physics already. This note is a friendly textbook-style
introduction to the four neural-network architectures we plan to train
for the FXT centrality paper. By the last page you should be able to
sit at a group meeting, hear someone say ``we ran a DeepSets model
against a Set Transformer with an evidential head,'' and translate that
into something concrete. We assume you have never built a neural
network. We do not assume you remember any ML jargon, and we will
define each term once when it first appears.

% =====================================================================
\section{The ML basics, in two pages}

\subsection{What ``training a neural network'' actually means}

A neural network is a parameterized function $f_\theta(x)$. The vector
$\theta$ is a list of numbers (the \emph{weights}, sometimes millions
of them) that we are free to choose. The function takes some input $x$
(in our case, features of a heavy-ion event) and returns a prediction
(in our case, the impact parameter $b$ and the centrality percentile).

The shape of $f_\theta$ is built out of two simple operations applied
in alternation: a linear transformation $y = Wx + c$, and a fixed
nonlinear ``activation'' applied element-wise (for example, ReLU,
defined as $\max(0, z)$). Stack a few dozen of these layers and the
overall function becomes flexible enough to approximate essentially
any smooth mapping. Importantly, the network is just arithmetic --- no
mystery, no hidden physics.

\emph{Training} is the process of choosing $\theta$ so that the
network's predictions match the truth on a set of events for which we
know the answer. We measure the mismatch with a \emph{loss function}
$L(\theta)$, for example the mean squared error
$L = \langle (b_{\text{pred}} - b_{\text{true}})^2 \rangle$. Then we
adjust $\theta$ in the direction that decreases $L$. The directional
derivative $\partial L / \partial \theta$ is the \emph{gradient}, and
the iterative recipe
\[
   \theta \leftarrow \theta - \eta \, \nabla_\theta L
\]
is \emph{gradient descent}. The number $\eta$ is the
\emph{learning rate}, typically $10^{-3}$ to $10^{-4}$. In practice we
average the gradient over a small batch of events (say 256) at a time;
this is \emph{stochastic gradient descent} and is essentially the
entire training algorithm.

To avoid the model memorizing the training set, we split our 350k
SMASH events into three disjoint pieces: \textbf{train} (80\,\%, used
to update weights), \textbf{validation} (10\,\%, used to pick
hyperparameters and decide when to stop), and \textbf{test} (10\,\%,
touched exactly once, at the very end, to report performance). Mixing
these up is the single most common methodological mistake in applied
ML, and we go to lengths to prevent it.

\subsection{Permutation invariance, in physics language}

A heavy-ion event is fundamentally a \emph{set} of final-state
particles. There is no first particle, second particle, third
particle. If I reshuffle the list, the event is the same event. Any
function we use to extract $b$ from the particle list had better have
the same property:
\[
   f(p_1, p_2, \dots, p_N) = f(p_{\pi(1)}, p_{\pi(2)}, \dots, p_{\pi(N)})
\]
for any permutation $\pi$. This is \emph{permutation invariance}, and
it is the central inductive bias in three of our four architectures.
``Inductive bias'' just means a structural assumption baked into the
model so it does not have to learn that assumption from data.

A vanilla MLP (multi-layer perceptron) given a flat vector of particle
4-momenta does \emph{not} have this property: swap particles 7 and 8,
and the inputs to all downstream neurons change. Three of our four
architectures fix this by construction. The fourth (the MLP baseline)
sidesteps it by feeding only permutation-invariant \emph{summary}
features (counts, means, sums).

A useful analogy: computing the total charge of a system from a list
of constituents is automatically permutation-invariant, because
summation is. Computing the total $\pT$ is the same. So is the mean
rapidity. These are all examples of the same universal recipe:
particle-wise function, then symmetric pool, then a function of the
pool. This is exactly the DeepSets template (Sec.~\ref{sec:deepsets}).

\subsection{What the four models share}

By design (see \texttt{project\_scope\_v2.md}, locked 2026-05-17), all
four architectures see the same inputs. This is so that the only
difference between them is the inductive bias, not the data.

\begin{itemize}[leftmargin=*]
\item \textbf{Per-particle features} (variable length, one row per
      charged hadron passing $0 \le \etalab \le 1.5$\footnote{The
      project scope memo records the wider window
      $0<\etalab<2$; the current code uses $0\le\etalab\le 1.5$ per
      STAR FXT TPC (Kimelman thesis); this discrepancy is to be
      reconciled.}):
      $\pT$, $\etalab$, $\phi$, charge.\footnote{Frame convention:
      SMASH and UrQMD generate events in collider (CM) mode with two
      opposing beams. For FXT-faithful analysis we Lorentz-boost into
      the target-rest frame using $y_{\text{boost}} =
      \mathrm{arccosh}(\sqsnn / 2 m_N)$ along the beam axis; $\etalab$
      denotes pseudorapidity in that boosted frame and matches the
      geometry STAR uses at fixed-target energies. The boost is
      implemented in \texttt{scripts/build\_padded\_cache.py}. (Note:
      a parallel CM-frame pipeline also exists in the codebase for
      cross-checks; the primer's quoted numbers are lab-frame unless
      stated otherwise.)}
      No PDG, no mass --- we deliberately exclude particle ID at the
      centrality stage, so the method does not depend on PID
      performance. At $\sqsnn=3.5$\,GeV the most central $5\,\%$ of
      events have a median of $\sim 150$ such tracks in this
      acceptance (measured from
      \texttt{data/processed/cached/all\_lab\_truth.h5}, top-5\,\% by
      truth $b$, $0\le\etalab\le 1.5$); a peripheral event can have
      zero.
\item \textbf{Event-level features} (always four scalars per event):
      $\sqsnn$, multiplicity in the acceptance window
      ($\text{mult}_{\text{lab}}$), mean $\pT$, and total $\pT$.
      $\sqsnn$ is what makes joint multi-energy training possible.
\end{itemize}

The MLP uses only the event-level features (plus a handful of derived
summaries). The other three architectures see the full particle list
\emph{and} the event-level features, the latter concatenated after the
permutation-invariant pool.

\subsection{What the four models output (the heads)}

All four models share the same two output heads:

\begin{enumerate}[leftmargin=*]
\item An \textbf{evidential regression head} for $b$. Instead of
      outputting a single number, it outputs four numbers
      $(\mu, \nu, \alpha, \beta)$ that parameterize a
      Normal-Inverse-Gamma (NIG) distribution. The practical
      consequence is that every prediction comes with its own
      uncertainty bar --- one number for the best-estimate
      $b$ and another for ``how sure the model is.'' This is what
      lets us draw a coverage plot later. (Caveat: a recent paper
      [Amini-NIG ``mirage,'' arXiv:2510.18322] argues NIG uncertainty
      can be miscalibrated; we will cross-check with deep ensembles
      for the headline result.)
\item A \textbf{centrality percentile head}. A softmax classifier over
      10 or 20 bins (0--5\,\%, 5--10\,\%, $\dots$). The softmax
      produces a probability for each bin, which we read either as a
      hard prediction (\textit{argmax}) or as a soft distribution.
\end{enumerate}

The total loss is a weighted sum of the regression loss (NIG negative
log-likelihood plus a regularizer) and the cross-entropy loss for the
classifier. We train both heads jointly.

\clearpage
% =====================================================================
\section{Architecture 1 -- MLP feature-vector baseline}
\label{sec:mlp}

\subsection{Input}

A flat vector of $\sim$10--20 scalar event-level summary features:
$\sqsnn$, charged multiplicity in $\etalab$ windows
(several sub-windows: 0--0.5, 0--1, 0--1.5), mean $\pT$, total $\pT$, the
ratio of forward to mid-rapidity multiplicity, the number of protons,
the number of antiprotons, and a small handful of similar hand-built
quantities. Crucially, \emph{no} particle-list input.

\subsection{Core idea}

The MLP --- multi-layer perceptron --- is the workhorse of plain-vanilla
deep learning. It is a stack of dense layers:
\[
   h_{\ell+1} = \sigma(W_\ell h_\ell + c_\ell)
\]
where $\sigma$ is a nonlinearity (we use ReLU), $W_\ell$ is a weight
matrix, $c_\ell$ a bias vector, and $h_\ell$ is the hidden
representation at layer $\ell$. We use 2--4 layers, widths of
128--512, and a small amount of \emph{dropout} (randomly zeroing 10\%
of neurons during training) as a regularizer.

There is no per-particle reasoning here at all. The model sees the
same RefMult observable a Glauber-RefMult fit uses, plus several
additional event-level summary scalars --- so strictly more
information than a single-observable Glauber fit. The inductive bias
is essentially zero --- we are just
giving the model a nice flexible parametric form that can learn the
mapping (summary numbers) $\to$ (b, percentile) better than a
hand-tuned threshold.

\subsection{Physicist analogy}

This is the ML equivalent of polynomial regression on engineered
features. If you have ever fit a yield versus centrality with a
cubic in mean $\pT$ and multiplicity, you have already done this in
spirit. The MLP just lets the polynomial be much higher order and
much more flexible.

\subsection{Why we picked it}

It is the lower-bound baseline. The MLP tells us: \emph{how much
performance comes from having any flexible nonlinear model at all,
vs. how much actually requires per-particle information?} If the MLP
on summary statistics already matches the more sophisticated
permutation-invariant models, then per-particle information is not
the key ingredient at FXT, and the paper's story shifts to
``calibrated UQ + multi-energy training do the work, even on summary
features.'' If the MLP loses, we have quantified the value of going
particle-level.

\subsection{Known limitations}

It can only ever exploit information present in the hand-built summary
features. Any signal living in the joint distribution of individual
particle kinematics (e.g.\ angular correlations, $\pT$ spectrum shape)
is invisible to it. At FXT energies, where the per-event statistics
are small (of order $10^2$ tracks even in central events) and every
track potentially matters, this is a real handicap.

% =====================================================================
\section{Architecture 2 -- DeepSets}
\label{sec:deepsets}

\subsection{Input}

The variable-length particle list directly:
$\{(\pT^{(i)}, \etalab^{(i)}, \phi^{(i)}, q^{(i)})\}_{i=1}^N$
where $N$ varies event by event. We pad shorter events to the same
length with masked-out zeros for efficient batching, but the network
ignores the padding.

\subsection{Core idea}

Zaheer et al.\ (NeurIPS 2017) proved a clean theorem: \emph{any}
permutation-invariant function of a set can be written in the form
\[
   f(\{x_1, \dots, x_N\}) = \rho\!\left(\sum_{i=1}^N \phi(x_i)\right)
\]
where $\phi$ and $\rho$ are arbitrary functions (in practice, both are
MLPs). The model first applies the \emph{same} small MLP $\phi$ to
each particle independently, producing a per-particle embedding (say,
a 64-dim vector). It then \emph{pools} these embeddings by summing
(or averaging, or max-pooling --- any symmetric operation) across the
list. The result is a single fixed-size event embedding, which is
then passed through another MLP $\rho$ to produce the output.

The sum step is what guarantees permutation invariance: $\sum_i$ does
not care about the order of the summands.

\subsection{Physicist analogy}

You have computed countless event-summed observables already. The
total transverse energy $\sum_i E_{\text{T}}^{(i)}$ is exactly a
DeepSets-style sum with $\phi(p) = E_{\text{T}}(p)$ and
$\rho(\cdot) = \mathrm{identity}$. The flow vectors
$Q_n = \sum_i e^{in\phi_i}$ are the same pattern with
$\phi(p) = e^{in\phi_p}$. DeepSets says: \emph{let the network choose
$\phi$ and $\rho$ for you, by gradient descent, optimized for the task
you actually care about}. Instead of hand-designing an event-summed
observable, you let the data tell you which 64 summed observables
matter for predicting $b$.

\subsection{Why we picked it}

It is the cleanest permutation-invariant architecture and the natural
``next step up'' from the MLP. The comparison MLP-vs-DeepSets isolates
the value of per-particle information \emph{cleanly}, because
DeepSets only differs from the MLP in having an actual per-particle
stage. It is also fast, easy to train, and easy to interpret --- you
can inspect what $\phi$ has learned.

\subsection{Known limitations}

The pool is a bottleneck. Two events that have the same per-particle
sum but different per-particle structure (for example, two narrow
back-to-back jets vs.\ an isotropic distribution with the same total
$\pT$) collapse to the same event embedding. The model has no way to
capture interactions \emph{between} particles, only their independent
contributions. For our purposes this may not matter much --- FXT
events are not jet-like --- but it is the theoretical ceiling on what
DeepSets can do.

\clearpage
% =====================================================================
\section{Architecture 3 -- Set Transformer}
\label{sec:settransformer}

\subsection{Input}

Same as DeepSets: the variable-length per-particle list.

\subsection{Core idea}

The Set Transformer (Lee et al., ICML 2019) replaces DeepSets'
independent per-particle embedding with one in which particles can
\emph{see each other} through attention. The headline operation is
the \emph{self-attention block}. For a list of $N$ particles, each
particle $i$ updates its embedding by taking a weighted average of
\emph{all other} particles' embeddings, where the weights are learned
based on which pairs are most relevant:
\[
   h_i^{\text{new}} = \sum_{j=1}^N \alpha_{ij}\, V(x_j),
   \qquad
   \alpha_{ij} = \frac{\exp(Q(x_i) \cdot K(x_j))}{\sum_k \exp(Q(x_i) \cdot K(x_k))}.
\]
The functions $Q$, $K$, $V$ (``query, key, value'') are learned linear
maps. After several such blocks, the per-particle embeddings encode
not just the particle's own properties but its relationship to every
other particle in the event. The model finally pools using
\emph{PMA} (pooling by multi-head attention), which is itself a
learned permutation-invariant operation.

The whole thing is still strictly permutation invariant, because
attention is computed by a sum that does not depend on particle order
--- only the values of $Q, K, V$ matter, which are particle-local
quantities.

\subsection{Physicist analogy}

Imagine you wanted to estimate centrality not just from each
particle's $\pT$ in isolation, but from the spatial pattern of
particles relative to each other (the gradient of multiplicity in
$\phi$, two-particle correlations, the presence of a recoil jet, the
azimuthal asymmetry that signals elliptic flow). Each of these
requires comparing particles against other particles. Attention is the
machinery that lets the network learn which pairs of particles matter
for the task.

The cost: self-attention scales as $\mathcal{O}(N^2)$ in particles per
event, which is fine for $N \sim 100$ FXT events but expensive for
$N \sim 1000+$.

\subsection{Why we picked it}

Two reasons. First, it is the strongest permutation-invariant
baseline in modern ML --- if particle-particle correlations carry real
centrality signal, this is the architecture that will find it.
Second, the comparison Set Transformer vs.\ DeepSets isolates the
value of \emph{interactions} beyond simple sums. If they tie,
volume-fluctuation signal at FXT lives in the marginals, not in
correlations. If the Set Transformer wins, correlations matter, and
that itself is a publishable physics finding.

\subsection{Known limitations}

More parameters, slower to train, slightly more sensitive to
hyperparameters than DeepSets. The interpretability is worse:
inspecting what attention has learned is a research topic in its own
right. There is also a small risk of overfitting on small per-event
$N$ because attention has high capacity relative to the data per
event.

% =====================================================================
\section{Architecture 4 -- Energy Flow Network (EFN)}
\label{sec:efn}

\subsection{Input}

Same per-particle list as DeepSets and Set Transformer, but EFN
distinguishes the ``energy'' (intensity) feature from the ``angle''
(direction) features. In our setup, $\pT$ plays the role of energy,
and $(\etalab, \phi, q)$ play the role of angles. (Charge isn't an
``angle'' geometrically, but functionally it behaves the same way: a
non-energy attribute of the particle.)

\subsection{Core idea}

The EFN (Komiske, Metodiev, Thaler --- JHEP 2019) is a
\emph{constrained} DeepSets. Its functional form is
\[
   \mathrm{EFN}(\{p_i\})
   = F\!\left(\sum_{i=1}^N z_i \, \Phi(\hat{n}_i)\right),
\]
where $z_i = \pT^{(i)}$ is the per-particle energy weight,
$\hat{n}_i$ are the angular variables, $\Phi$ is a small MLP from
angles to a latent vector, and $F$ is a larger MLP from the pooled
vector to the output.

Compare to the DeepSets form, $\rho(\sum_i \phi(p_i))$: the EFN takes
$\phi(p_i) = z_i \, \Phi(\hat{n}_i)$, i.e.\ the per-particle embedding
must be linear in the energy. That linearity constraint is the
key. It makes the network \emph{IRC-safe}: infrared-and-collinear
safe. In jet-physics language, IRC safety means the output is
continuous under (i) the addition of a zero-energy particle
(infrared safety) and (ii) the exact collinear splitting of one
particle into two particles sharing the same total energy (collinear
safety). Strictly, IRC safety guarantees only these limits --- it
does \emph{not} by itself guarantee robustness to a finite
$\pT > 50\,\text{MeV}$ cut or to finite-width detector smearing.
Separately from IRC safety, the EFN output is a continuous (smooth
MLP) function of the per-particle $\pT$ values, so small numerical
perturbations in $\pT$ produce small changes in output; this is the
property we actually exploit in the Phase B detector-emulation
comparison.

\subsection{Physicist analogy}

The most natural example is the standard jet observable
$\sum_i \pT^{(i)} f(\eta_i, \phi_i)$ --- this is exactly EFN form,
with the function $f$ replaced by a learnable MLP $\Phi$. If you have
used jet substructure observables like $N$-subjettiness, you have used
an IRC-safe construction in spirit. EFN is the deep learning version:
a function in this safe family, optimized end-to-end.

\subsection{Why we picked it}

Three reasons.

\begin{enumerate}[leftmargin=*]
\item It is the standard permutation-invariant architecture in
      experimental HEP. Reviewers (especially LHC-side
      machine-learners) will ask ``did you benchmark against EFN'' ---
      we want a yes.
\item IRC safety is a useful property for Phase B of our project, when
      detector emulation kicks in. EFN by construction should degrade
      gracefully under the $\sim 1\,\%$ $\pT$ smearing and tracking
      inefficiency we apply. If it does (and DeepSets/Set Transformer
      do not), that is informative.
\item It positions us against the CBM PointNet line. PointNet is
      \emph{not} IRC safe; using EFN explicitly contrasts our
      methodology with theirs.
\end{enumerate}

\subsection{Known limitations}

The linearity-in-energy constraint costs some flexibility. If the
optimal solution requires nonlinear weighting of particle energies
(say, particles below a threshold should be ignored entirely), EFN
cannot represent that exactly. Komiske et al.\ also introduced
\emph{PFN}, an unconstrained variant; we use the constrained EFN
because the constraint is the whole point. Also, the choice of
``energy'' weight $z_i$ is itself a modelling decision: the original
EFN paper uses $z_i = \pT^{(i)}/\pT^{\text{jet}}$ or
$E_i/E^{\text{jet}}$, and we use unnormalized $z_i = \pT^{(i)}$.
That is a valid choice but it does mean the IRC-safety statement is
relative to ``soft in $\pT$,'' not ``soft in total energy.'' At FXT
energies this distinction matters less than at the LHC because the
particles are non-relativistic enough that $\pT$, $E$, and $\pT/M$
ordering of particles is similar; we still expect the structural
inductive bias to be useful.

\clearpage
% =====================================================================
\section{Side-by-side comparison}
\label{sec:comparison}

\begin{table}[h]
\centering
\renewcommand{\arraystretch}{1.35}
\begin{tabularx}{\textwidth}{l|X|X|X|X}
\toprule
              & \textbf{MLP}
              & \textbf{DeepSets}
              & \textbf{Set Transformer}
              & \textbf{EFN} \\
\midrule
\textbf{Input type}
              & Scalar summary features ($\sim 10$--20 numbers)
              & Variable-length particle list
              & Variable-length particle list
              & Particle list, energy + angle split \\
\midrule
\textbf{Inductive bias}
              & None beyond smoothness of the function
              & Permutation invariance via symmetric sum/mean pool
              & Permutation invariance + learned particle--particle interactions
              & Permutation invariance + linearity in energy (IRC safe) \\
\midrule
\textbf{Rough parameter count}
              & $10^4$--$10^5$
              & $10^5$--$10^6$
              & $10^6$--$10^7$
              & $10^5$--$10^6$ \\
\midrule
\textbf{Expected strength at FXT}
              & Trivially fast; transparent; lower-bound baseline
              & Best simplicity-to-power ratio; clean comparison to MLP
              & Captures correlations if they matter (flow, jet-like patterns)
              & Robust to detector smearing; HEP-standard \\
\midrule
\textbf{Expected weakness at FXT}
              & Blind to per-particle structure; can only learn from hand-engineered summaries
              & Pool is a bottleneck; no particle--particle interactions
              & Most parameters; risk of overfitting on small $N$; harder to interpret
              & Linearity-in-energy constraint may cost flexibility; IRC safety less natural at FXT than for jets \\
\bottomrule
\end{tabularx}
\caption{The four architectures at a glance. ``Rough parameter
count'' is the order of magnitude we expect at the configurations we
plan to train; final values depend on width and depth, which are
hyperparameters.}
\end{table}

% =====================================================================
\section{What we expect to learn from running all four}

Each pairwise comparison answers a distinct physics-or-method
question.

\begin{itemize}[leftmargin=*]
\item \textbf{MLP vs.\ DeepSets} --- does per-particle information
      beat hand-engineered summary features at FXT? At top RHIC and
      LHC the answer is unambiguously yes. At low multiplicity it is
      genuinely open: a 100-track central event has very few moments
      that are not already captured by mean, sum, and a couple of
      $\eta$-window multiplicities.

\item \textbf{DeepSets vs.\ Set Transformer} --- do particle--particle
      correlations carry centrality signal beyond the marginals? If
      the Set Transformer wins by a noticeable margin, that is itself
      a finding worth a paragraph in the discussion section, because
      it tells the cumulant-analysis community something about what
      they are throwing away when they collapse to multiplicity.

\item \textbf{DeepSets vs.\ EFN} --- does the IRC-safety constraint
      cost us performance at truth level, and does it pay off under
      detector emulation in Phase B? The Phase A vs.\ Phase B
      comparison for EFN is one of the cleanest experiments in the
      project.

\item \textbf{All four vs.\ Glauber/RefMult baselines} --- this is the
      headline plot. We expect every ML model to beat the classical
      baselines in the peripheral regime where $\sim 47\,\%$ of FXT
      events have zero charged particles in $|\eta_{\text{CM}}|<0.5$
      (see CLAUDE.md; SMASH dumps CM-frame 4-momenta, so this is the
      CM-frame mid-rapidity acceptance). The interesting question is
      how the four
      stratify among themselves.

\item \textbf{Calibrated UQ} --- across the four architectures, are
      the evidential uncertainty bars actually well-calibrated? If the
      Amini-NIG ``mirage'' critique bites, we will see it as
      systematic over- or under-coverage, and we will fall back to
      deep ensembles for the headline numbers.
\end{itemize}

The four-architecture comparison is not just hedging: it is a
controlled experiment in which the inputs and the heads are held
fixed, and only the inductive bias varies. That is what makes the
result publishable as a methodology study rather than as a one-off
benchmark.

% =====================================================================
\section*{Further reading}

\begin{itemize}[leftmargin=*, itemsep=2pt]
\item Zaheer et al., \emph{Deep Sets}, NeurIPS 2017,
      \texttt{arXiv:1703.06114}.
\item Lee et al., \emph{Set Transformer}, ICML 2019,
      \texttt{arXiv:1810.00825}.
\item Komiske, Metodiev, Thaler, \emph{Energy Flow Networks: Deep
      Sets for Particle Jets}, JHEP 01 (2019) 121,
      \texttt{arXiv:1810.05165}.
\item Amini et al., \emph{Deep Evidential Regression}, NeurIPS 2020,
      \texttt{arXiv:1910.02600}.
\item The Amini-NIG critique: \texttt{arXiv:2510.18322} ---
      background reading on why we back up evidential UQ with deep
      ensembles.
\item Mallick, Tripathy, Behera, et al., \emph{Estimation of impact
      parameter and transverse spherocity in heavy-ion collisions
      using machine learning}, \texttt{arXiv:2103.01736} --- the
      closest prior art at top RHIC/LHC.
\item Omana Kuttan et al., \emph{A fast centrality-meter for the CBM
      experiment}, Phys.\ Lett.\ B 811 (2020) 135872,
      \texttt{arXiv:2009.01584} --- the closest prior art in the FXT
      regime, using PointNet.
\end{itemize}

\end{document}
